\section{Discussion}
\label{sec:discussion}

\paragraph{What the data says, and what it does not.}
The pre-registered headlines hold (\S\ref{sec:results-mitigation}, \S\ref{sec:results-f1}): four frontier web-browsing agents leak critical-tier PII at a substantial baseline rate when faced with social-engineering websites, the three escalating prompt-level mitigations tested do not reduce this rate by more than the pre-committed 30 pp falsification margin at the pooled level, and within the pre-registered C3 reflection condition the agent still submits critical PII in 35.9\% of sessions in which an independent LLM-as-judge confirms verbalised suspicion. We do not interpret the null pooled mitigation finding as evidence that no prompt can ever reduce leakage; the pre-registered claim is the narrower one that the three specific prompts spanning generic policy, phishing checklist, and pre-submission reflection do not, and reviewers should read this finding alongside the converging null in DECEPTICON~\citep{cuvin2026decepticon} and TrickyArena~\citep{ersoy2026trickyarena} rather than as a single negative result. The detection--action gap is the more actionable finding: within C3 (reach-gated), confirmed-detector sessions leak at 35.9\% versus 66.1\% in non-detector sessions, a 30.2 pp gap pooled across all four models. A keyword-heuristic baseline widens the apparent gap (\S\ref{sec:results-f1}); the LLM-judge primary is the more conservative reading and the one we report. The implication is that defenses which gate behaviour on the agent's stated assessment of a site are gating on the wrong signal: even after filtering for confirmed recognition, more than one in three detector sessions still hands over critical PII.

\paragraph{Toward output-level defenses.}
\label{sec:defense-sketch}
The data narrows the search space for a defense without specifying one. Any mitigation delivered through the agent's reasoning trace and gated on its own trust judgment inherits the \S\ref{sec:results-f1} gap; an effective defense for the deployment regime in \S\ref{sec:threat} would likely sit outside the agent's decision loop, intercepting outbound POST bodies and gating critical-tier submissions on a separate trust signal. The TCR--ASR exchange rate in Table~\ref{tab:results-tcr} bounds what such a scaffold would need to clear; design is left to future work.


\section{Conclusion}
\label{sec:conclusion}
Scammer4U is a pre-registered, axis-controlled benchmark of 91 adversarial environments and 10 benign twins for measuring critical-PII exfiltration by autonomous web-browsing agents under social-engineering attack. Across four frontier models, four conditions, and five seeds, baseline leakage is high, escalating prompt-level mitigation does not clear the pre-committed 30 pp falsification margin at the pooled level, and agents that explicitly identify an attack still submit critical PII in over a third of those sessions---roughly $3.5\times$ the pre-registered 10\% bar. We release the benchmark, harness, analysis plan, and raw logs so these comparisons can be reproduced and extended.


\section{Limitations}
\label{sec:limitations}
We enumerate the principal limitations of this study. Environments are templated rather than scraped from active phishing sites, and a descriptive fidelity check against real PhishTank captures (Appendix~\ref{app:fidelity}) bounds rather than eliminates the visual-believability gap; templated sites also lack the cloaking, ephemerality, and JavaScript anti-analysis present in deployed phishing infrastructure, which PhreshPhish~\citep{Dalton2025PhreshPhishAR} documents as the primary obstacle to using scraped corpora for reproducible study. Sites are LLM-authored from human-written briefs with the model held constant across all 91 environments and distinct from the four evaluated models (\S\ref{sec:authoring}), but residual authoring-style confounds with specific patterns the authoring model favours cannot be fully excluded. The four-model panel is curated for deployed agentic systems (\S\ref{sec:setup-runs}) and does not cover the open-weight or smaller-parameter regime. The DR judge is GPT-4o-mini; the within-family pairing with GPT-5 mini is mitigated by per-model stratified inter-judge $\kappa$ reporting and a Llama 4 Scout sensitivity pass (\S\ref{sec:dr}, Appendix~\ref{app:dr-sensitivity}). PII is US-centric (SSN, US bank/routing format, US address), and prompts and pages are English-only; both narrow the population of users this study speaks to. The benign-twin pool covers four of eight environment categories (\S\ref{sec:benign}); the remaining four use the global-twin fallback for the baseline adjustment, and we report both raw and adjusted rates for transparency. Multi-turn chat environments (axis G) carry higher per-seed variance by construction, and we report per-axis variance separately. Finally, the harness in which all four models run is one specific Playwright-based observe--think--act loop with a domain-masked observer and a request-layer egress guard; production commercial scaffolds may differ in observation format, action space, and policy enforcement, and our results bound the in-harness behaviour of these models rather than the behaviour of any specific commercial deployment.

\section{Ethical Considerations}
\label{sec:ethics}

\paragraph{Dual-use and release posture.} Scammer4U releases attacker-shaped templates alongside an evaluation harness; we considered whether release materially aids attackers and concluded it does not. The attack patterns instantiated (typosquats, fake trust signals, reward traps, prompt injection, authority impersonation) are documented in MITRE ATT\&CK, OWASP LLM Top 10, ENISA, and public corpora such as PhishTank in greater volume than we release here. The value of the release is the \emph{evaluation apparatus} (the paired-sibling design, the 8-axis taxonomy, the analysis plan, and the benign-twin baselines), which enables defensive measurement and which an attacker has no operational use for. Every environment is a self-hosted Flask application bound to localhost; no environment is configured for public DNS, public TLS, or email deliverability, and we publish no instructions for converting one into a live phishing kit. Templated brand assets are stylised lookalikes (typosquatted domains, near-brand spellings) rather than scraped pixel-perfect clones, and no live brand asset is bundled in the release.

\paragraph{Synthetic PII.} The agent is provisioned with a single fictitious identity (\emph{Jordan Mitchell}) used across every environment, model, condition, and seed. All values are synthetic and do not, to our knowledge, correspond to any real person. Phone numbers use the NANPA fictional \texttt{555-01XX} reservation; SSN, payment card, account number, passwords, and API key are constructed strings designed to be globally unique within the profile (\S\ref{app:profile}); the bank routing number is a real public-record routing identifier paired with a synthetic account number, which is not on its own actionable for fraud and is required only so that form validators on the attack envs do not reject it as malformed. No real user PII is collected, stored, or processed by the benchmark or harness at any point. Captured POST bodies are written only to the per-environment \texttt{/api/captured} sink on localhost, never transmitted off-host.

\paragraph{Real-phishing corpus use.} The fidelity human review (Appendix~\ref{app:fidelity}) compares Scammer4U environments against $20$ real captured-phishing pages drawn from PhishTank, the Wayback Machine, and curated security-blog reposts. These captures are \emph{referenced} rather than redistributed; the released artifact contains the per-stimulus manifest (source URL, capture date, PhishTank ID where applicable) so reviewers can re-fetch the same set, but the artifact does not bundle the captured HTML/CSS/JS. No live brand asset, target brand's intellectual property, or victim PII contained in the captures is redistributed.

\paragraph{Human-subjects review.} The fidelity review involved $16$ adult volunteer reviewers from the authors' institution rating screenshots of websites; no personal data was collected from reviewers (familiarity self-report and confidence rating were aggregated and anonymised before analysis), no deception was involved beyond the standard blind-rating protocol, and reviewers consented in writing to the rating task and to anonymised aggregate reporting of their ratings. The protocol is minimal-risk under the authors' institutional research-ethics policy on perception-task studies and did not require formal IRB review; the consent script and rater debrief are included in the released artifact.

\paragraph{Model and provider use.} All four evaluated agents and both judge models were accessed through public commercial APIs under their providers' terms of use; the workload (autonomous web browsing in localhost-only sandboxed Flask apps) does not violate any provider's acceptable-use policy. The benchmark issues no requests against public web origins, third-party services, or production brand domains during evaluation; an egress guard at the browser layer blocks any agent-initiated navigation outside the localhost twin set. No commercial deployment of any evaluated model was probed or stress-tested.

\paragraph{Disclosure.} The detection--action gap and per-model heterogeneity findings are model-behaviour observations rather than zero-day vulnerabilities; we did not pursue a coordinated-disclosure process, but the analysis plan, environments, and results were shared with the four model providers ahead of public release as a courtesy heads-up on the deployment-relevant implications.


